\chapter{Модели эффективности доступа к многомерным данным}
\label{ch:models}

\section{Модель стоимости агрегирующего запроса}
% Развитие модели кандидатской: число просмотренных узлов как функция
% размерности D, селективности, ветвистости f и перекрытия ключей.

\section{Ветвистость узла и её связь с размером ключа}
% f = f(размер страницы, размер ключа, служебные представления ключа).
% Отсюда — задача главы 3.

\section{Модель стоимости построения индекса}
% Последовательная вставка против сортировки: O(N log_f N) обращений к страницам
% против стоимости внешней сортировки. Отсюда — задача главы 4.

\section{Модель стоимости сопровождения}
% Обход в логическом порядке против физического; стоимость случайного чтения
% против последовательного; влияние упреждающего чтения. Отсюда — задача главы 5.

\section{Выводы по главе}
% Три параметра, на которые можно воздействовать без изменения класса
% операторов: ветвистость, перекрытие, порядок обращения к страницам.
% Каждой — по главе.
